% =====================================================================
\chapter{Introduction}
\label{ch:introduction}
% =====================================================================

\section{Background and Context}
\label{sec:intro-background}

Modern organisations process a continuous stream of \emph{visually-rich
documents}---forms, invoices, receipts, contracts, and reports---whose meaning
depends not only on the words they contain but on \emph{where} those words sit on
the page. Extracting structured information from such documents is the task of
\emph{key-information extraction} (KIE), the cornerstone of document-understanding
systems deployed in finance, logistics, healthcare, and public administration.
The dominant approach casts KIE as token-level sequence labelling over a tri-modal
input: text obtained by optical character recognition (OCR), the two-dimensional
layout encoded as bounding boxes, and the page image. Multimodal document
encoders such as LayoutLMv3~\citep{huang2022layoutlmv3}, LiLT~\citep{wang2022lilt},
and BROS~\citep{hong2022bros} fuse these signals and define the state of the art.

Deployed document systems are not static. New document types appear, annotation
schemas are extended, and the distribution of incoming documents drifts as an
organisation enters new markets or adopts new templates. A model must therefore
\emph{evolve}---learning new entity types and new domains over time---without
re-training from scratch on all historical data, which is often impossible
because earlier data is unavailable, expensive to store, or restricted for
privacy reasons. This is the setting of \emph{continual learning} (CL): learning a
sequence of tasks while retaining performance on those seen earlier.

The central obstacle to continual learning is \emph{catastrophic
forgetting}~\citep{mccloskey1989catastrophic}: when a neural network is fine-tuned
on a new task, gradient updates overwrite the parameters that encoded earlier
tasks, and performance on those tasks collapses. A large literature has developed
remedies---regularising important weights~\citep{kirkpatrick2017ewc},
distilling from past models~\citep{li2017lwf}, replaying stored
exemplars~\citep{rolnick2019er,buzzega2020der}, learning task-specific
prompts~\citep{wang2022l2p,wang2022dualprompt,smith2023coda}, or allocating
low-rank adapters per task~\citep{wang2023olora}. Almost all of these methods,
however, were developed and evaluated on \emph{unimodal} image- or text-classification
benchmarks, and they treat the network as a homogeneous block: regularisation is
applied uniformly across parameters, or layers are frozen by hand-tuned heuristics.

For a \emph{tri-modal} document encoder this uniform treatment is unsatisfying,
because the network is manifestly heterogeneous: it contains a text stream, a
visual/patch stream, layout (2D position) embeddings, and cross-modal fusion
layers that combine them. It is not known which of these components forgets
most, nor whether the pattern is stable across continual scenarios. Answering
that question is a prerequisite for designing a continual-learning method that
spends its capacity where forgetting actually happens rather than spreading it
thinly across the whole network. This thesis takes a \emph{measure-first,
then-fix} stance: it first \emph{diagnoses} where forgetting lives in a multimodal
document encoder, and only then proposes a remedy whose design is justified by
the diagnosis.

\section{Objectives}
\label{sec:intro-objectives}

The thesis pursues the following concrete objectives:

\begin{enumerate}
  \item \textbf{Formulate continual document KIE as a unified problem.} Cast
        key-information extraction across heterogeneous datasets (FUNSD, CORD,
        SROIE) as continual token-level sequence labelling, and define three
        continual scenarios---class-incremental (CIL), domain-incremental (DIL),
        and a mixed scenario---over a shared output format
        (Chapter~\ref{ch:experimental-setup}).
  \item \textbf{Build a per-component forgetting diagnostic.} Develop an
        instrumentation protocol that attributes forgetting to individual
        components of a tri-modal encoder using representational drift (linear
        CKA), per-parameter-group Fisher information, and standard forgetting
        metrics, under a controlled modality-ablation design
        (Chapter~\ref{ch:methodology}, \S\ref{sec:method-diagnostic}).
  \item \textbf{Characterise where forgetting concentrates.} Apply the
        diagnostic to naive sequential fine-tuning and test whether forgetting is
        uniform across components or concentrates in specific ones, establishing
        the stability of any finding across seeds and task orders
        (Chapter~\ref{ch:results}, \S\ref{sec:results-diagnosis}).
  \item \textbf{Derive and evaluate a mechanism-targeted remedy.} Propose a
        continual-learning method whose mechanism is concentrated on the
        dominant-forgetting component identified by the diagnosis, and evaluate it
        against ten baselines spanning the major CL families
        (Chapters~\ref{ch:methodology} and~\ref{ch:results}).
  \item \textbf{Establish a reproducible benchmark.} Release the FCS-CL benchmark
        (FUNSD/CORD/SROIE continual scenarios), the diagnostic tooling, and all
        configurations, with results reported as means over three seeds with
        statistical tests (Chapters~\ref{ch:implementation}
        and~\ref{ch:experimental-setup}).
\end{enumerate}

\section{Achievements}
\label{sec:intro-achievements}

% Numeric claims marked TODO are filled once experiments complete (Chapter 6).

\noindent\textbf{A per-component forgetting diagnosis for tri-modal document
encoders.} We introduce, to our knowledge, the first protocol that decomposes
catastrophic forgetting in a multimodal document encoder into its text, visual,
layout, and fusion contributions. By masking input modalities of a \emph{single}
LayoutLMv3 architecture we isolate each modality's contribution without the
confound of comparing different architectures.
% TODO: state headline finding + effect size once pilot is analysed.

\noindent\textbf{A mechanism-targeted continual-learning method (DocCL).} Guided
by the diagnosis, we design a method that concentrates its continual-learning
mechanism on the dominant-forgetting component rather than treating the network
uniformly.
% TODO: state the improvement over the best component-agnostic baseline
% (e.g. "+X entity-F1, -Y BWT on 2/3 scenarios") once the main grid is complete.

\noindent\textbf{The FCS-CL benchmark and reproducible tooling.} We provide a
unified continual-learning benchmark over three public document datasets with
three scenarios, ten baselines, and per-component instrumentation, all under a
single configuration system and tracked experiment log, enabling like-for-like
comparison that the fragmented prior literature has lacked.

\section{Overview of the Dissertation}
\label{sec:intro-overview}

\textbf{Chapter~\ref{ch:background}} reviews document understanding and
information extraction, the continual-learning problem and its method families,
the nascent literature on continual learning for information extraction, evidence
on where forgetting lives in multimodal models, and the backbones relevant to
this work; it closes by identifying the research gaps the thesis addresses.
\textbf{Chapter~\ref{ch:methodology}} presents the two methodological
contributions: the per-component forgetting diagnostic and the mechanism-targeted
continual-learning method derived from it.
\textbf{Chapter~\ref{ch:implementation}} describes the framework, the
document-IE pipeline, and the system architecture, including the instrumentation
hooks and experiment tracking.
\textbf{Chapter~\ref{ch:experimental-setup}} details the datasets, the continual
task construction, the model and training configuration, the baselines, and the
evaluation methodology.
\textbf{Chapter~\ref{ch:results}} reports the diagnosis as the headline result,
compares the remedy against baselines, presents ablations and a computational-cost
analysis, and discusses implications and threats to validity.
\textbf{Chapter~\ref{ch:conclusions}} states limitations, conclusions, and
future work.
